\documentclass[11pt]{article}
\usepackage[margin=1in]{geometry}
\usepackage{lmodern,microtype,enumitem}
\title{Q11 Discussion Questions}
\date{}
\begin{document}
\maketitle
\begin{enumerate}[noitemsep]
\item Kalai and Vempala show calibrated models \emph{must} hallucinate on arbitrary facts, yet post-trained models hallucinate less and are measurably \emph{less} calibrated. Is trading calibration for caution the right deal? What does a user lose when a model's probabilities stop meaning what they say?
\item The module argues hallucination is not lying because there is no intent. Where does that line actually sit? If a model reliably maintains a false claim under clarification (``are you sure?'') because refusals are low-probability under its distribution, is that closer to confabulation, stubbornness, or deception --- and does the moral category change any engineering decision?
\item Retrieval can inject errors: Mallen et al.\ found retrieval \emph{hurts} on popular entities the model already knows, and a RAG system inherits its index's failures. Given that, when would you \emph{not} ground a production system --- and what audit would you run on the index before trusting it?
\item Gekhman et al.\ find fine-tuning on new factual knowledge linearly increases hallucination. If SFT teaches the form of answers rather than facts, what should a domain-adaptation team do whose corpus is 90\% new knowledge (e.g., internal company facts)? Defend a concrete pipeline.
\item The module's taxonomy separates fabrication, misattribution, and faithful errors. Design an evaluation that reports all three separately for a RAG assistant --- what labeled data would you need, and which of TruthfulQA, cloze probes, or HaluEval could you reuse for which slice?
\end{enumerate}
\end{document}
